% Original-Goodman contribution to planned164.9; not a substitute for Schwartz's full 2011 chapter.
\chapter{Goodman, Projectibility and the Choice of Predicates}
\label{sec:p15-goodman-projectibility}

An observation does not carry its own unique description into an unobserved
case. The same inspected emerald can be described with predicates whose
extensions agree on the inspected objects and disagree elsewhere. Goodman's
new riddle asks what distinguishes the descriptions that may legitimately
be projected. Counting favourable instances, requiring a universal sentence,
or removing an overt date from its spelling does not by itself answer that
question. The issue concerns which generalizations are candidates for
confirmation, not only how much confirmation a chosen candidate receives
\parencite[III, \S4]{GoodmanNewRiddleTeachingScan}.

The primary presentation used here is chapter~III, sections~1--5, of
\emph{Fact, Fiction and Forecast}. The available chapter scan does not identify
its edition; a reference to a 1963 publication prevents treating its filename
as evidence of an unchanged 1955 text. This unidentified scan does not
establish the chronology of the book's earliest editions.
The fourth-edition publisher's cover dates the book's first publication
to 1954 \parencite[back cover]{Goodman1983FourthProjection}.
Goodman's argument, its independently stated formal examples below, and
Schwartz's later historical discussion have different evidential roles.
The examples prove the specified logical and probabilistic statements;
they do not prove a normative criterion of projectibility.

\section{The problem a projectibility account must address}
\label{sec:p15-goodman-task}

\begin{hbdefinition}{def:p15-goodman-projection}{Projection and the three assessment questions}
A \emph{projection} extends a description of some given cases to cases not
given by that evidence. Its target need not lie in the future: another place,
an unexamined contemporary object or an unobserved earlier occurrence may
be the target. For a hypothesis and evidence, distinguish:

\begin{enumerate}
\item whether the hypothesis is \emph{projectible}, an eligible basis for the
intended extension;
\item whether, and to what degree, the specified evidence supports it;
\item whether the projected claim is true.
\end{enumerate}

Goodman's \emph{lawlike} versus \emph{accidental} distinction addresses the
first question. In his account lawlikeness is not synonymous with truth,
importance, or possession of a universal grammatical form.
\end{hbdefinition}

Goodman's response to the older problem compares the adjustment of
inference rules with the adjustment of accepted inferences. It is an account
of justification through their mutual accommodation, not a promise of
infallible anticipation. This philosophical proposal is distinct from
Hume's original argument in
\cref{sec:p15-hume-and-the-problem-of-induction}; it must not be silently
substituted for Hume's own treatment. Nor does the proposal say that every
habit or every currently accepted inference is valid. Disagreement about
which judgments to retain remains possible
\parencite[III, \S\S1--2]{GoodmanNewRiddleTeachingScan}.

For a concrete procedure, the inputs therefore have to include more than
a list of successes. They include the candidate vocabulary, the background
claims, the target cases, and whatever preference among competing
projections is being defended. Specifying these inputs makes a procedure
well defined; it does not establish that its preferences are justified.

\section{Grue has an observation condition, not a colour-change instruction}
\label{sec:p15-goodman-grue}

\begin{hbdefinition}{def:p15-goodman-grue}{Fixed-cutoff observation predicates}
Fix an object domain $U$ and a cutoff $t$. Write $\mathsf{O}_t(x)$ for
``$x$ was examined before $t$,'' $\mathsf{E}(x)$ for emerald,
$\mathsf{G}(x)$ for green and $\mathsf{B}(x)$ for blue. Define
\[
 \mathsf{Q}_t(x)\ \Longleftrightarrow\
 (\mathsf{O}_t(x)\land\mathsf{G}(x))\lor
 (\neg\mathsf{O}_t(x)\land\mathsf{B}(x)).
\]
Here $\mathsf{Q}_t$ represents \emph{grue}. The cutoff and the observation
history are fixed in this comparison. An object unexamined before $t$
belongs to the second branch even if it existed long before $t$.
The definition does not instruct an object to change its colour when a
clock reaches $t$.

The two universal candidates are
\[
 H_G:\ \forall x(\mathsf{E}(x)\to\mathsf{G}(x)),\qquad
 H_Q:\ \forall x(\mathsf{E}(x)\to\mathsf{Q}_t(x)).
\]
When incompatibility is claimed, assume explicitly that no emerald is
both green and blue. No assumption that these exhaust all colours is needed.
\end{hbdefinition}

\begin{hbtheorem}{thm:p15-goodman-grue-models}{proved=annex:p15-goodman-grue-models}{Same inspected instances, opposed extensions}
Let $D$ be a nonempty finite set of inspected emeralds, and $b\notin D$
an emerald not examined before $t$. There are two structures on
$U=D\cup\{b\}$ satisfying exactly the same positive observations
$\mathsf{E}(a),\mathsf{O}_t(a),\mathsf{G}(a)$ for every $a\in D$,
with $\mathsf{O}_t$ having extension $D$, such that one satisfies $H_G$
and refutes $H_Q$, and the other satisfies $H_Q$ and refutes $H_G$.
In both structures every recorded emerald is also grue. More generally,
under green--blue incompatibility and the existence of an emerald not
examined before $t$, $H_G$ and $H_Q$ cannot both hold.
\end{hbtheorem}

This construction identifies the missing choice in an unrestricted
positive-instance recipe: if it treats all substituted predicates alike,
each recorded instance qualifies for both universal candidates. It does
not establish equal posterior probabilities, equal predictive accuracy,
or equal rational support under a specified background theory. Those are
additional claims, requiring additional assumptions. If there are no
unexamined emeralds, the displayed universal candidates need not disagree.
Likewise a statement that \emph{all} previously examined objects are green
is not entailed by observations confined to emeralds.

\begin{hbexample}{def:p15-goodman-antecedent}{Changing the projected class as well as its colour description}
Let $\mathsf{V}(x)$ mean rose and define
\[
 \mathsf{A}_t(x)\ \Longleftrightarrow\
 (\mathsf{O}_t(x)\land\mathsf{E}(x))\lor
 (\neg\mathsf{O}_t(x)\land\mathsf{V}(x)).
\]
All the recorded emeralds in the preceding construction are
$\mathsf{A}_t$ and $\mathsf{Q}_t$. A recipe permitting the generalization
$\forall x(\mathsf{A}_t(x)\to\mathsf{Q}_t(x))$ from those instances
therefore predicts blue for a rose not examined before $t$: the rose
satisfies the second branch of $\mathsf{A}_t$, and its grueness reduces
to blueness. The premises contain no observed rose. This is the
antecedent counterpart of the grue problem, corresponding to Goodman's
\emph{emerose} example, not an empirical claim about roses
\parencite[III, \S4, n.~10]{GoodmanNewRiddleTeachingScan}.
\end{hbexample}

\section{Why a syntactic or qualitative vocabulary test is not neutral}
\label{sec:p15-goodman-language}

\begin{hbdefinition}{def:p15-goodman-bleen}{A common-domain grue--bleen pair}
On the same domain and with the same observation condition, set
\[
 \mathsf{R}_t(x)\ \Longleftrightarrow\
 (\mathsf{O}_t(x)\land\mathsf{B}(x))\lor
 (\neg\mathsf{O}_t(x)\land\mathsf{G}(x)).
\]
This is the \emph{bleen} counterpart. The common-domain definition is
an explicit formal comparison: the source introduces grue for things
generally and illustrates the reverse vocabulary on emeralds.
\end{hbdefinition}

\begin{hblemma}{thm:p15-goodman-vocabulary}{proved=annex:p15-goodman-vocabulary}{Invertibility of the colour descriptions}
For arbitrary extensions of $\mathsf{G},\mathsf{B},\mathsf{O}_t$,
\[
 \begin{split}
 \mathsf{G}(x)&\ \Longleftrightarrow\
 (\mathsf{O}_t(x)\land\mathsf{Q}_t(x))\lor
 (\neg\mathsf{O}_t(x)\land\mathsf{R}_t(x)),\\
 \mathsf{B}(x)&\ \Longleftrightarrow\
 (\mathsf{O}_t(x)\land\mathsf{R}_t(x))\lor
 (\neg\mathsf{O}_t(x)\land\mathsf{Q}_t(x)).
 \end{split}
\]
Thus either pair may describe the other relative to the observation
predicate. The identities require neither colour exhaustiveness nor
green--blue incompatibility. Moreover, in first-order logic with equality,
any sentence $H$ is equivalent to $H\land(c=c)$ for a newly introduced
individual name $c$.
\end{hblemma}

A rule inspecting only whether the displayed sentence contains an
individual or temporal name is therefore representation-sensitive. A
primitive predicate may conceal a defining condition which mentions a
cutoff; a logically equivalent sentence may gratuitously display an
individual name. Declaring green and blue primitive while grue and bleen
are defined distinguishes the spellings, but reversing that declaration
reverses the test. The lemma does not show that green and blue lack a
physical distinction, or that all semantically informed restrictions
must fail. It shows exactly why a restriction based on this choice of
primitive definitions needs a justification beyond its syntax
\parencite[III, \S4]{GoodmanNewRiddleTeachingScan}.

Invoking an \emph{overhypothesis} raises the same dependency at another
level. Grouping observations by material, rather than by occupancy of
a particular desk, already selects a grouping to project. A higher-order
classification may supply useful background; merely calling that
classification lawlike does not explain why it is entitled to that role.

\section{Confirmation rules and total evidence}
\label{sec:p15-goodman-confirmation}

\begin{hbtheorem}{thm:p15-goodman-confirmation-collapse}{proved=annex:p15-goodman-confirmation-collapse}{Two unrestricted closure demands trivialize confirmation}
Let $C(e,h)$ be a relation between formulas of classical logic. Suppose
it satisfies both conditions, without restrictions on $e,h,k$:
\begin{enumerate}
\item if $h\models e$, then $C(e,h)$;
\item if $C(e,h)$ and $h\models k$, then $C(e,k)$.
\end{enumerate}
Then $C(e,h)$ holds for every pair of formulas $e,h$.
\end{hbtheorem}

The result concerns the conjunction of the two specified demands, not
every confirmation theory and not every restricted consequence principle.
Goodman's discussion uses it to expose a defective definition. Qualifying
the instance universe and taking the total stated evidence into account
are different repairs, not a proof that all such repairs are sufficient
\parencite[III, \S3]{GoodmanNewRiddleTeachingScan}.

For example, one object described as black and another as nonblack do not
together instantiate a universe in which everything is black. Projecting
separately from each isolated report discards the opposing report.
Similarly, a nonblack nonraven logically satisfies the contrapositive
of the black-raven universal, but that equivalence alone assigns no
amount of evidential support. Background assumptions about possible
ravens and the observation protocol matter to such an assignment.

\begin{hbexample}{def:p15-goodman-support}{Identical sample agreement need not mean identical evidential support}
Consider the three possible structures $w_G,w_Q,w_N$. The first two are
the structures of \cref{thm:p15-goodman-grue-models}; the third has a
recorded emerald which is neither green nor blue. Let the evidence event
$d$ be that all recorded emeralds are green. Then
$d=\{w_G,w_Q\}$, $H_G=\{w_G\}$ and $H_Q=\{w_Q\}$.
Put prior probabilities $3/5,3/10,1/10$ on these three structures.
Direct conditioning gives
\[
 P(d)=9/10,\quad P(H_G\mid d)=2/3,\quad P(H_Q\mid d)=1/3.
\]
The increases are respectively $1/15$ and $1/30$. Both candidates fit
the whole recorded sample, both have likelihood one for that sample,
and both increase in probability, but their posterior probabilities
and absolute increases differ. They have the same multiplicative
increase, $10/9$; even the choice of support measure must be stated.
These are consequences of this declared finite model, not an assertion
that these priors are the uniquely rational emerald priors.
\end{hbexample}

\begin{hbwarning}{def:p15-goodman-extension}{Sample agreement, full extension and evidence are different relations}
For predicates $F,K$, equality of
$\{a\in D:F(a)\}$ and $\{a\in D:K(a)\}$ is agreement on $D$.
Equality of their extensions in $U$ is a stronger, model-relative claim.
Logical equivalence requires agreement in every admissible structure.
If two hypothesis events are identical in a fixed probability space,
their conditional probabilities are identical whenever defined. None
of these equalities follows merely from a finite sample on which the
predicates agree. \Cref{thm:p15-goodman-grue-models} supplies the counterexample
to the full-extension inference; the probability example supplies the
counterexample to an inference of equal support from sample fit.
\end{hbwarning}

\section{Entrenchment, conflicts and the remaining normative choice}
\label{sec:p15-goodman-entrenchment}

Schwartz's historical discussion identifies \emph{entrenchment} with the
role of predicates in earlier projections, rather than with the mere
number of times a word is uttered. It reports inheritance through more
general predicates and through predicates taken to be coextensive.
Consequently entrenchment is not reducible to the present emerald sample,
nor to a guarantee that an entrenched predicate will predict successfully
\parencite[\S6, p.~402]{Schwartz2011GoodmanProjectibility}.

This matters when a replacement predicate has the same extension.
If its full coextension with an entrenched predicate is established,
an account may transfer the relevant projection history. Agreement only
on the inspected emeralds does not establish the premise for that
transfer: grue and green in the two constructed structures agree on $D$
but not on $U$. Confusing these two kinds of agreement would erase
exactly the contested extension.

The same discussion reports a revision associated with Schwartz,
Scheffler and Goodman in 1970. In its condensed account, conflicting
hypotheses are compared by entrenchment; a better-entrenched competitor
overrides its rival. A hypothesis with all rivals overridden is classified
as projectible, an overridden one as unprojectible, and an unresolved
conflict can make a hypothesis non-projectible. These are reported
historical classifications, not three grades of posterior probability.
The condensed account explicitly omits further complications. It must
not be represented as the complete original chapter-IV construction,
or as an executable decision procedure proved adequate for arbitrary
languages \parencite[\S6, p.~402 and n.~12]{Schwartz2011GoodmanProjectibility}.

An application must actually specify which hypotheses conflict, their
applicable projection histories, the comparison of entrenchment, and
the treatment of unresolved comparisons. Calling a pair ``compatible''
or one predicate ``natural'' does not perform those tasks. Conversely,
encoding those choices in a vocabulary, prior or model class does not
make them disappear: it makes them explicit premises. A statistical
or learning guarantee can remain mathematically correct relative to
its stated model while not establishing that the selected vocabulary
is the projectible one.

\begin{proofboundary}{def:p15-goodman-normative-boundary}{What the constructions establish}
The full proofs here establish finite underdetermination of extensions,
invertible descriptions and the collapse of one unrestricted confirmation
relation. The finite probability calculation is conditional on its priors
and space of possibilities. None proves that green is normatively
projectible, that entrenchment is a complete justification, or that every
semantic criterion fails. The full revised fourth-edition entrenchment
construction is developed in
\cref{sec:p15-goodman-entrenchment-and-projectibility}. The original
first- and second-edition rules, the distinct 1970 paper and the complete
Schwartz historical treatment require their own source definitions and
arguments; the bounded historical report above is not a replacement for them.
\end{proofboundary}
